% Template from latex.tugraz.at

\documentclass[%
%handout, % prints handouts (=no animations, for printed version)
%mathserif
%xcolor=pst,
14pt
% fleqn
]{beamer}

\usepackage{beamerthemedefault}

\useoutertheme[subsection=false]{smoothbars}
\useinnertheme[shadow=true]{rounded}

\setbeamercolor{block title}{fg=black,bg=gray}
\setbeamercolor{block title alerted}{use=alerted text,fg=black,bg=alerted text.fg!75!bg}
\setbeamercolor{block title example}{use=example text,fg=black,bg=example text.fg!75!bg}
\setbeamercolor{block body}{parent=normal text,use=block title,bg=block title.bg!25!bg}
\setbeamercolor{block body alerted}{parent=normal text,use=block title alerted,bg=block title alerted.bg!25!bg}
\setbeamercolor{block body example}{parent=normal text,use=block title example,bg=block title example.bg!25!bg}

%% for including video/audio
%\usepackage{multimedia} %needs hyperref-package

%% aktive Referenzen
\usepackage{hyperref}

%\usepackage{ngerman}			% language set to new-german
\usepackage[english]{babel}
\usepackage[utf8]{inputenc} 	% coding of german special characters
% \usepackage{ae,aecompl}
% \usepackage{amsmath,amssymb,amstext} 	% support for mathematics

% \usepackage{listings}		% include programming code
%\usepackage{amsfonts}		% blackboard fonts: $\mathbb{N,Z,R,C,...}$
%\usepackage{latexsym}
%\usepackage{textcomp}
%\usepackage{mathptmx,courier}
			% \textdegree \textcelsius \textperthousand
			% \copyright \texttrademark \textregistered
			% \textmu (non-italic mu)
%\usepackage{geometry}	% change paper dimension an margins
%    \geometry{verbose,paperwidth=128mm,paperheight=90mm}
%    \geometry{tmargin=0mm,bmargin=0mm,lmargin=10mm,rmargin=10mm}
%\usepackage{graphicx}  % \includegraphics[options]{file.eps}
			% options = scale, width, totalheight, height, depth,
			% angle = deg, origin = {l c r}{top, Baseline, bottom}
			% \scalebox{h-scale}[v-scale]{object}
			% \rotatebox[x=xdim,y=ydim]{angleCCW}[object}
%\usepackage{tabularx}	% \begin{tabular}{...X...} stretches column
%\usepackage{multirow}
%\usepackage{floatflt}
%\usepackage{hhline}
%\usepackage{colortbl}
\usepackage{array}
\usepackage{setspace}

%\usepackage[thinspace,thinqspace,squaren,textstyle]{SIunits}

\definecolor{tug}{rgb}{0.96862,0.14509,0.27450}

\setbeamertemplate{headline}[text line]{
	\begin{beamercolorbox}[wd=\paperwidth,ht=8ex,dp=4ex]{}
		\insertnavigation{0.85\paperwidth} 
		\raisebox{-10pt}{\includegraphics[width=15mm]{tuglogo}}\vskip2pt
		\hskip-1pt\rule{\paperwidth}{0.3pt}
	\end{beamercolorbox}
}

\setbeamertemplate{navigation symbols}{}

\definecolor{gray}{rgb}{0.8,0.8,0.8}
\setbeamercolor{footline}{fg=black,bg=gray}

% Fußzeile mit Autor, Titel und Foliennummer / Gesamtfolienzahl
\setbeamertemplate{footline}[text line]{
	\hskip-1pt
	\begin{beamercolorbox}[wd=\paperwidth]{footline}
			\rule{\paperwidth}{0.3pt}
			\colorbox{tug}{\rule{3pt}{0pt}\rule{0pt}{3pt}}
			\textbf{\rule{0pt}{5pt}\insertshortauthor\hfill\insertshortinstitute\hfill%
					\insertshorttitle\rule{1em}{0pt}}
			\rule{\paperwidth}{0.3pt}
	\end{beamercolorbox}
	\begin{beamercolorbox}[wd=\paperwidth,ht=2ex,dp=2ex]{white}
	\end{beamercolorbox}
}%

%% Titelblatt-Einstellungen
\title{Implementation of an Evolutionary Optimization Framework and Application to the Open Job Shop Problem}
%\subtitle{}
\author{Joris Bayer, Stefan Kroboth}
% \date{}		% wenn ein anderes als das heutige Datum eingesetzt werden soll

% Subject und Keywords für PDF-Datei
\subject{Implementation of an Evolutionary Optimization Framework and Application to the Open Job Shop Problem}
\keywords{Numerical Optimization, Evolutionary Algorithms, Open Job Shop Problem}

\titlegraphic{\includegraphics[width=20mm]{tuglogo}}

%%%%%%%%%%%%%%%%%%%%%%%%%%%%%%%%%%%%%%%%%%%%%%%%%%%%%%%%%%%%%%%%%%%%%%%%%%%%
\begin{document}
%%%%%%%%%%%%%%%%%%%%%%%%%%%%%%%%%%%%%%%%%%%%%%%%%%%%%%%%%%%%%%%%%%%%%%%%%%%%

\begin{frame}[plain]
  \frametitle{}
  \titlepage % erzeugt Titelseite
\end{frame}



\begin{frame}
  \frametitle{Content}
        \tableofcontents[hideallsubsections %
                        % ,pausesections
                        ] % erzeugt Inhaltsverzeichnis
\end{frame}

\section{Introduction}
\subsection{nix}
\begin{frame}
  \frametitle{Motivation}
\end{frame}

\begin{frame}
  \frametitle{Language}
\end{frame}

\section{evolib}
\subsection{nix}
\begin{frame}
  \frametitle{evolib - A Julia Library}
\end{frame}

\begin{frame}
  \frametitle{Algorithms}
\end{frame}

\begin{frame}[fragile]
  \frametitle{Example}
  %\begin{exampleblock}{Beispiel-Block}
  \vspace{-0.8cm}
{\scriptsize
\begin{verbatim}
load("evolib.jl")

function rosenbrock(chr::Chromosome)
    x1 = chr[1].gene
    x2 = chr[2].gene
    chr.fitness = 100*(x2-x1^2)^2+(1-x1)^2-50*((x1+1)^2+(x2-1)^2)
                 -10*((x1-1.5)^2+(x2-2.5)^2)+exp(2*x2-5)
end

function rosenbrock(pop::Population)
    for i=1:length(pop)
        rosenbrock(pop[i])
    end
end

popul = rand(Population, 100, 2, rosenbrock, 1.0, 6.0, -6.0 )

@time evo_slash(popul, 2, 150, 1000, 0.85, 0.0000001, rosenbrock)

\end{verbatim}}
  %\end{exampleblock}
\end{frame}

\section{Open Job Shop Problem}
\subsection{Introduction}
\begin{frame}
  \frametitle{Problem definition}
\end{frame}

\subsection{Algorithms}
\begin{frame}
  \frametitle{Permutation Genetic Algorithm}
\end{frame}

\begin{frame}
  \frametitle{Hybrid Genetic Algorithm}
\end{frame}

\begin{frame}
  \frametitle{Selfish Gene Algorithm}
\end{frame}

\begin{frame}
  \frametitle{Other approaches}
\end{frame}

\section{Evaluation}
\subsection{nix}
\begin{frame}
  \frametitle{Evaluation}
\end{frame}

\begin{frame}
  \frametitle{Random Number Generator}
\end{frame}

\begin{frame}
  \frametitle{Convergence}
\end{frame}

\begin{frame}
  \frametitle{Comparison}
\end{frame}

\section{Conclusion}
\subsection{nix}
\begin{frame}
  \frametitle{Conclusion}
\end{frame}

%%%%%%%%%%%%%%%%%%%%%%%%%%%%%%%%%%%%%%%%%%%%%%%%%%%%%%%%%%%%%%%%%%%%%%%%%%%%
%\section{Beispiele}
%%%%%%%%%%%%%%%%%%%%%%%%%%%%%%%%%%%%%%%%%%%%%%%%%%%%%%%%%%%%%%%%%%%%%%%%%%%%%

%\subsection{Aufzählungen}
%\begin{frame}
	%\frametitle{Aufzählungen erstellen}
	%\begin{spacing}{1.5}
	%\begin{itemize}
		%\item Punkt 1
		%\begin{itemize}
			%\item Unterpunkt 1
			%\item Unterpunkt 2
		%\end{itemize}
		%\item Punkt 2
		%\item Punkt 3
	%\end{itemize}
	%\end{spacing}
%\end{frame}

%%%%%%%%%%%%%%%%%%%%%%%%%%%%%%%%%%%%%%%%%%%%%%%%%%%%%%%%%%%%%%%%%%%%%%%%%%%%%

%\subsection{Blöcke}
%\begin{frame}
	%\frametitle{Blöcke}
	%\framesubtitle{Normale Blöcke}
	%\begin{block}{Normaler Block}
		%Das ist ein Block mit einer Formel
		%\begin{eqnarray*}
		%u(x,t) & = & \sum_{k=1}^{\infty} f_k \sin \frac{k \pi x}{L} \cos 
					 %\frac{k \pi t}{aL} + \\
				 %& + & \sum_{k=1}^{\infty} g_k \sin \frac{k \pi x}{L} \sin 
					 %\frac{k \pi t}{aL} \\
		%\end{eqnarray*}
	%\end{block}
%\end{frame}

%%%%%%%%%%%%%%%%%%%%%%%%%%%%%%%%%%%%%%%%%%%%%%%%%%%%%%%%%%%%%%%%%%%%%%%%%%%%%

%\begin{frame}[fragile]
	%\frametitle{Blöcke}
	%\framesubtitle{Beispiel-Blöcke}
	%\begin{exampleblock}{Beispiel-Block}
	%\begin{semiverbatim}
%SUCHE (A,x)
%1: i = 0
%2: WHILE i<n
%3:     i = i+1
%4:     \alert{IF A[i]=x THEN RETURN i}
%5: ELSE RETURN -1
	%\end{semiverbatim}
	%\end{exampleblock}
%\end{frame}

%%%%%%%%%%%%%%%%%%%%%%%%%%%%%%%%%%%%%%%%%%%%%%%%%%%%%%%%%%%%%%%%%%%%%%%%%%%%%

%\begin{frame}
	%\frametitle{Blöcke}
	%\framesubtitle{Alert-Blöcke}
	%\Large
	%\begin{alertblock}{Alert-Block}
		%\centering
		%Achtung! Nicht zu viel Inhalt auf eine Folie!
	%\end{alertblock}
%\end{frame}

%%%%%%%%%%%%%%%%%%%%%%%%%%%%%%%%%%%%%%%%%%%%%%%%%%%%%%%%%%%%%%%%%%%%%%%%%%%%%

%\subsection{Spalten und Grafiken}
%\begin{frame}
	%\frametitle{Spalten und Grafiken}
	%\begin{columns}
		%\begin{column}{0.5\textwidth}
			%Aufzählung in Spalte 1
			%\begin{enumerate}
				%\item Punkt 1
				%\item Punkt 2
				%\item Punkt 3
			%\end{enumerate}
		%\end{column}
		%\begin{column}{0.5\textwidth}
			%\begin{center}
			%%\includegraphics[width=0.4\textwidth]{smail}

			%Grafik in Spalte 2
			%\end{center}
		%\end{column}
	%\end{columns}
%\end{frame}

%%%%%%%%%%%%%%%%%%%%%%%%%%%%%%%%%%%%%%%%%%%%%%%%%%%%%%%%%%%%%%%%%%%%%%%%%%%%%
%\section{Abschnitt 2}
%%%%%%%%%%%%%%%%%%%%%%%%%%%%%%%%%%%%%%%%%%%%%%%%%%%%%%%%%%%%%%%%%%%%%%%%%%%%%

%%%%%%%%%%%%%%%%%%%%%%%%%%%%%%%%%%%%%%%%%%%%%%%%%%%%%%%%%%%%%%%%%%%%%%%%%%%%%
%\section{Abschnitt 3}
%%%%%%%%%%%%%%%%%%%%%%%%%%%%%%%%%%%%%%%%%%%%%%%%%%%%%%%%%%%%%%%%%%%%%%%%%%%%%

%%%%%%%%%%%%%%%%%%%%%%%%%%%%%%%%%%%%%%%%%%%%%%%%%%%%%%%%%%%%%%%%%%%%%%%%%%%%%
%\section{Zusammenfassung}
%%%%%%%%%%%%%%%%%%%%%%%%%%%%%%%%%%%%%%%%%%%%%%%%%%%%%%%%%%%%%%%%%%%%%%%%%%%%%

%\begin{frame}
	%\frametitle{Zusammenfassung}
	%\begin{itemize}
		%\item Punkt 1
		%\item Punkt 2
		%\item Punkt 3
	%\end{itemize}
%\end{frame}


%%%%%%%%%%%%%%%%%%%%%%%%%%%%%%%%%%%%%%%%%%%%%%%%%%%%%%%%%%%%%%%%%%%%%%%%%%%%%
%\section*{} % Damit die Kontakt-Folie nicht im Inhaltsverzeichnis aufscheint
%%%%%%%%%%%%%%%%%%%%%%%%%%%%%%%%%%%%%%%%%%%%%%%%%%%%%%%%%%%%%%%%%%%%%%%%%%%%%

%\begin{frame}
%\frametitle{Kontakt}
%nothing
%\end{frame}

%%%%%%%%%%%%%%%%%%%%%%%%%%%%%%%%%%%%%%%%%%%%%%%%%%%%%%%%%%%%%%%%%%%%%%%%%%%%
\end{document}
%%%%%%%%%%%%%%%%%%%%%%%%%%%%%%%%%%%%%%%%%%%%%%%%%%%%%%%%%%%%%%%%%%%%%%%%%%%%

%% EOF
